\chapter{Annotazioni e Convenzioni}
\label{app:A}

Questa appendice raccoglie le convenzioni usate nella descrizione del sistema
RAG. L'obiettivo è evitare ripetizioni nei capitoli principali e rendere
esplicito il significato di nomi, formati e parametri ricorrenti.

\section{Repository e componenti}

La Tabella~\ref{tab:repo_conventions} fissa i nomi dei repository usati nei
capitoli principali.

\begin{table}[ht]
  \centering
  \caption{Convenzioni di denominazione dei repository.}
  \label{tab:repo_conventions}
  \small
  \begin{tabularx}{\textwidth}{L{4.8cm}Y}
    \toprule
    \textbf{Nome} & \textbf{Significato} \\
    \midrule
    \texttt{i3guild}
      & Applicazione web principale, basata su Next.js, PostgreSQL e Keycloak. \\
    \breakcode{i3guild-embedding-pipeline}
      & Job Python che estrae, normalizza e indicizza il corpus delle gilde. \\
    \texttt{i3guild-rag}
      & Microservizio FastAPI che espone retrieval ibrido, chat RAG e
        generazione controllata. \\
    \bottomrule
  \end{tabularx}
\end{table}

\section{Formato dei documenti indicizzati}

I documenti salvati in Qdrant derivano dai record applicativi di i3Guild. Ogni
documento contiene un testo normalizzato e un insieme di metadati usati per
filtri, diagnostica e deduplicazione.
La Tabella~\ref{tab:document_metadata_conventions} elenca i campi principali
richiamati nella descrizione del payload Qdrant.

\begin{table}[ht]
  \centering
  \caption{Metadati principali associati ai documenti indicizzati.}
  \label{tab:document_metadata_conventions}
  \small
  \begin{tabularx}{\textwidth}{L{4cm}Y}
    \toprule
    \textbf{Campo} & \textbf{Uso} \\
    \midrule
    \texttt{guild\_id}
      & Identificativo della Gilda sorgente. \\
    \texttt{guild\_name}
      & Nome della Gilda, usato anche nei risultati mostrati all'utente. \\
    \texttt{epoch\_id}
      & Identificativo dell'epoca a cui appartiene la Gilda. \\
    \texttt{participants}
      & Lista dei partecipanti, utile come metadato filtrabile. \\
    \texttt{participation\_mode}
      & Modalità di partecipazione, ad esempio in presenza, remota o ibrida. \\
    \texttt{is\_individual}
      & Indica se la proposta riguarda una Gilda individuale. \\
    \texttt{chunk\_index}
      & Posizione del chunk quando una Gilda viene suddivisa. \\
    \texttt{chunk\_count}
      & Numero totale di chunk generati per la stessa Gilda. \\
    \bottomrule
  \end{tabularx}
\end{table}

\section{Eventi NDJSON}

Gli endpoint streaming restituiscono una sequenza di righe JSON indipendenti.
La scelta del formato NDJSON permette al frontend di elaborare subito i metadati
RAG e poi aggiornare l'interfaccia man mano che arrivano i token.
La Tabella~\ref{tab:ndjson_event_conventions} riassume gli eventi usati per
questa comunicazione incrementale.

\begin{table}[ht]
  \centering
  \caption{Tipi di evento usati dagli endpoint streaming.}
  \label{tab:ndjson_event_conventions}
  \small
  \begin{tabularx}{\textwidth}{L{3.5cm}Y}
    \toprule
    \textbf{Evento} & \textbf{Significato} \\
    \midrule
    \texttt{rag}
      & Contiene documenti recuperati, filtri, query riscritta e motivo del
        taglio dei risultati. \\
    \texttt{status}
      & Indica l'avanzamento di una generazione asincrona lato applicazione. \\
    \texttt{stream\_start}
      & Segnala l'inizio della risposta generata. \\
    \texttt{token}
      & Contiene un frammento incrementale della risposta. \\
    \texttt{stream\_end}
      & Chiude lo stream dopo l'ultimo token generato. \\
    \texttt{error}
      & Comunica un errore gestibile senza interrompere il contratto dello
        stream. \\
    \bottomrule
  \end{tabularx}
\end{table}
